% Non-Disclosure Agreement. Content transcribed from NDA.docx, which is
% the source of truth. Every party detail, date and place comes from
% registry/ and the person record -- nothing here is hard-coded.
\documentclass{paridyum}
\begin{document}

\paridyumletterhead
\paridyumdateref{\dkDocDate}{\dkDocRef}
\paridyumtitle{Non-Disclosure Agreement (NDA)}

This Non-Disclosure Agreement (``Agreement'') is entered into between
\textbf{\dkOrgName} (``Paridyum'') and \textbf{\dkPIdentityLegalName} (``Intern'').

\section*{1. Purpose}
During the internship, the Intern may receive access to private or confidential
information relating to Paridyum. This Agreement explains how that information
must be protected and used.

\section*{2. Confidential Information}
``Confidential Information'' means information relating to Paridyum that is not
publicly available. This may include:
\begin{itemize}[leftmargin=*, noitemsep]
    \item source code and software;
    \item product ideas and features;
    \item technical architecture;
    \item algorithms and technical research;
    \item designs and prototypes;
    \item internal documentation;
    \item business plans and strategies;
    \item financial information;
    \item customer or user information;
    \item marketing plans;
    \item internal communications;
    \item passwords, API keys, tokens, and credentials;
    \item security information;
    \item databases and datasets;
    \item unreleased products or services; and
    \item other information that would reasonably be considered confidential.
\end{itemize}
Information does not need to be specifically marked ``Confidential'' to be
protected under this Agreement.

\section*{3. Intern's Responsibilities}
The Intern agrees to:
\begin{itemize}[leftmargin=*, noitemsep]
    \item use Confidential Information only for authorized internship work;
    \item keep Confidential Information private;
    \item not share it with unauthorized persons;
    \item not publish it online;
    \item not upload private source code or documents to public repositories;
    \item not copy or store Confidential Information unnecessarily;
    \item protect passwords, API keys, and other credentials; and
    \item take reasonable steps to prevent unauthorized access or disclosure.
\end{itemize}

\section*{4. Information That Is Not Confidential}
Confidential Information does not include information that:
\begin{itemize}[leftmargin=*, noitemsep]
    \item was already publicly available without a breach of this Agreement;
    \item was lawfully known to the Intern before receiving it from Paridyum;
    \item was lawfully received from another person without a confidentiality
          obligation; or
    \item was independently developed by the Intern without using Paridyum's
          Confidential Information.
\end{itemize}

\section*{5. Third-Party and AI Services}
The Intern must not upload or provide Paridyum's Confidential Information to
unauthorized third-party services. This includes external AI tools, file-sharing
services, code-sharing platforms, or similar services unless their use has been
approved by Paridyum.

\section*{6. Security Incidents}
The Intern must promptly inform Paridyum if they become aware of:
\begin{itemize}[leftmargin=*, noitemsep]
    \item leaked credentials;
    \item accidentally published source code;
    \item unauthorized access;
    \item lost or stolen devices containing Company information;
    \item accidental disclosure of confidential information; or
    \item any other suspected security incident.
\end{itemize}
Reporting an accidental incident promptly helps Paridyum take appropriate action
to limit potential harm.

\section*{7. Required Disclosure}
If the Intern is legally required by a court, government authority, or applicable
law to disclose Confidential Information, the Intern may disclose only the
information legally required. Where legally permitted, the Intern should notify
Paridyum before making the disclosure.

\section*{8. Return or Deletion of Information}
When the internship ends, or when requested by Paridyum, the Intern must:
\begin{itemize}[leftmargin=*, noitemsep]
    \item return Company documents and materials;
    \item stop accessing Company systems;
    \item remove Company credentials from personal devices;
    \item return Company equipment; and
    \item delete Confidential Information stored personally where requested.
\end{itemize}
The Intern may retain information where retention is legally required, provided
that the information continues to be protected.

\section*{9. Intellectual Property}
This NDA protects confidential information but does not itself determine ownership
of work created during the internship. Intellectual-property ownership is governed
by the separate \textbf{Intellectual Property Assignment Agreement}.

\section*{10. Duration}
The Intern's confidentiality obligations continue after the internship ends. For
ordinary Confidential Information, these obligations will continue for
\textbf{5 years after the internship ends}. Information that qualifies as a trade
secret or remains legally protected as confidential must continue to be protected
for as long as it retains that status.

\section*{11. Future Paridyum Entity}
If Paridyum later operates through a registered company or other legal entity, the
rights and protections under this Agreement may be transferred to that entity where
permitted by applicable law.

\section*{12. Governing Law}
This Agreement will be governed by the laws of \dkOrgJurisdictionCountry. Subject to
applicable law, disputes relating to this Agreement will be subject to the
appropriate courts in \textbf{\dkOrgJurisdictionForum}.

\section*{13. Acceptance}
By signing below, the Intern confirms that they have read and understood this
Agreement and agree to protect Paridyum's Confidential Information in accordance
with its terms.

\vspace{1.5em}
\paridyumofficers{\dkSigCeoName}{\dkSigCeoDesignation}{\dkSigCtoName}{\dkSigCtoDesignation}

\vspace{1.5em}
\paridyumcountersign{Intern}{\dkPIdentityLegalName}

\end{document}
